\chapter{Concluzii}
\label{ch:concluzii}
\pagestyle{fancy}

Lucrarea de față a pornit de la dorința de a realiza ceva concret și spectaculos, o
mână robotică ce reproduce în timp real mișcările mâinii umane, și a parcurs întregul
drum de la idee până la un sistem funcțional. Pe parcurs, proiectul a îmbinat domenii
foarte diferite, de la mecanica acționării prin tendoane și fabricarea prin printare 3D,
până la proiectarea digitală pe FPGA, filtrarea Kalman și electronica de izolare și
alimentare. Fiecare dintre aceste etape a venit cu propriile provocări, iar rezolvarea
lor, pas cu pas, a transformat un ansamblu de modele și componente disparate într-un
sistem unitar care funcționează.

\section{Aportul personal al lucrării}
\label{sec:concluzii-contributii}

Contribuția centrală a acestei lucrări o reprezintă partea de control implementată pe
placa FPGA, în care unghiurile primite de la sistemul de viziune sunt filtrate și
transformate în comenzi pentru servomotoare. În jurul acestei părți a fost dezvoltat un
filtru Kalman cu două stări, implementat integral în aritmetică cu virgulă fixă Q16.16,
capabil să prelucreze toate cele opt canale în interiorul unui singur ciclu de control.
Adaptarea modelului teoretic al filtrului la constrângerile unei platforme fără unitate
în virgulă mobilă, precum și încadrarea întregii logici în constrângerile temporale ale
unui circuit de 100~MHz, constituie principalul aport tehnic al lucrării.

Pe lângă filtru, a fost proiectat un întreg lanț de prelucrare, de la recepția și
validarea cadrelor prin UART, mecanismul de sincronizare de tip cutie poștală și
generarea celor opt semnale PWM, până la bucla de monitorizare prin convertorul
analog-digital. Închiderea temporală a circuitului, obținută prin descompunerea
operațiilor complexe, prin realizarea împărțirii pe mai multe cicluri și prin
distribuirea stării filtrului în memorii separate, a reprezentat una dintre cele mai
solicitante etape și, totodată, una dintre cele mai instructive.

O contribuție importantă este și realizarea fizică a mâinii, un ansamblu hibrid care
combină cele mai potrivite elemente din mai multe proiecte open-source, completate cu
piese proprii, modelate special pentru a face legătura între componente. Montarea
sistemului de acționare prin fire și arcuri, precum și realizarea unui sistem de
feedback prin intervenția directă asupra servomotoarelor, au necesitat un volum
considerabil de muncă practică. Nu în ultimul rând, integrarea componentelor Analog
Devices, atât pentru izolarea galvanică, cât și pentru achiziția și alimentarea de
precizie, a adăugat proiectului o dimensiune concretă de proiectare electronică.

\section{Evaluarea sistemului realizat}
\label{sec:concluzii-analiza}

Rezultatele obținute confirmă că sistemul își atinge scopul propus. Partea digitală
funcționează corect și stabil pe placa FPGA, respectând constrângerile temporale, iar
mâna reproduce mișcările comandate și reușește să prindă și să susțină obiecte reale.
Un rezultat îmbucurător a fost faptul că integrarea părții mecanice cu cea digitală a
funcționat de la prima încercare, ceea ce validează abordarea de testare pe niveluri,
în care fiecare subsistem a fost verificat înainte de a fi adus în ansamblu.

Există, desigur, și limitări. Prinderea este fermă pentru obiectele cu suprafețe care
oferă aderență, dar mai puțin sigură pentru obiectele netede din plastic lucios, care
tind să alunece, deoarece degetele nu dispun de un material antiderapant la vârfuri. De
asemenea, servomotoarele folosite, fiind modele accesibile, prezintă neuniformități care
au impus calibrări suplimentare, atât la nivelul semnalului PWM, cât și la nivelul
feedback-ului. Aceste limitări nu afectează funcționarea de ansamblu, dar indică
direcțiile în care sistemul poate fi îmbunătățit.

\section{Direcții viitoare}
\label{sec:concluzii-dezvoltare}

O primă direcție de dezvoltare vizează îmbunătățirea prinderii, prin adăugarea unui
material antiderapant la vârful degetelor, care ar permite manipularea sigură a unei
game mai largi de obiecte, inclusiv a celor netede. Tot pe partea mecanică, alimentarea
servomotoarelor la o tensiune mai ridicată, de 6~V, posibilă prin ajustarea modulului de
alimentare, ar crește cuplul disponibil și forța de strângere. O îmbunătățire și mai
consistentă ar fi înlocuirea servomotoarelor accesibile actuale cu modele de calitate
superioară, cu un cuplu mai mare și cu o cursă unghiulară mai amplă, care ar oferi atât
o forță de prindere sporită, cât și o libertate de mișcare mai apropiată de cea a mâinii
umane.

O altă direcție firească este eliminarea legăturii prin cablu dintre calculator și placă,
prin transmiterea datelor pe cale radio. Un modul de comunicație fără fir, conectat în
locul legăturii seriale actuale, ar permite operarea mâinii de la distanță, sporind
mobilitatea sistemului și apropiindu-l de scenariile reale de teleoperare.

Pe partea de control, filtrul Kalman ar putea fi extins la un model cu trei stări, care
să includă și accelerația, oferind o urmărire mai fidelă a mișcărilor rapide. Marja
foarte largă de resurse rămase pe circuitul FPGA permite cu ușurință o astfel de
extindere, fără a compromite respectarea constrângerilor de timp.

În final, cea mai firească direcție de continuare o reprezintă integrarea completă cu
partea de viziune, astfel încât mâna să urmărească în timp real și în mod robust
gesturile unui operator. Probele preliminare au arătat că lanțul întreg, de la captura
video până la mișcarea mâinii, poate funcționa unitar, iar rafinarea acestei integrări ar
transforma prototipul actual într-un sistem de teleoperare pe deplin funcțional,
apropiindu-l de scopul care a inspirat această lucrare.
